\documentclass[11pt]{article}
\usepackage{fontspec}
\usepackage[french,english]{babel}
\usepackage{amsmath}
\usepackage{amssymb}
\usepackage{amsthm}
\usepackage{mathrsfs}
\usepackage{tikz-cd}
\usepackage{enumitem}

\newcommand{\printedpage}[1]{\marginpar{\footnotesize p.~#1}}
\newcommand{\msp}[1]{\marginpar{\footnotesize [#1]}}
\newcommand{\Hotz}{\mathrm{Hot}^{\wedge}_{0}}
\newcommand{\Ch}{\underline{\mathrm{Ch}}}
\newcommand{\Ty}{\underline{\mathrm{Ty}}}
\newcommand{\CNo}{C_{\mathbf{N}^{o}}}
\newcommand{\lab}[2][l]{\text{\footnotesize\begin{tabular}{@{}#1@{}}#2\end{tabular}}}

\begin{document}

\printedpage{69}
% head: A. Grothendieck : Dérivateurs, Ch. VII, Édition des Universités de Montpellier II et Paris VII

\msp{page 108}
I find no direct arrow, for that matter, between $C$ and either of the objects $C_{\infty}$, $H_{C}$; all I see in this direction is
\[
\begin{tikzcd}[column sep=large, row sep=huge]
\CNo \arrow[r, "{\substack{\text{localization}\\ \text{functor}}}"] \arrow[d] & C_{\infty} \\
C, &
\end{tikzcd}
\]
where the vertical arrow is the projection arrow of $\CNo = C \times_{T} \mathbf{N}^{o}$ (with $\mathbf{N}^{o} \xrightarrow{\ t\ } T$). This establishes a close topological relation between $C$ (as an object of Cat\,hot or simply of Hot) and $C_{\infty}$, only if the two preceding arrows are at least Hot-equivalences. And the typical (problematic) case where the question of the relation between $C$ and $C_{\infty}$ arises, is therefore the one where $T = \mathbf{N}^{o}$, i.e. that of an arbitrary direct system of (small) categories, indexed by $\mathbf{N}$, and interpreted as defined by a fibered category $C = \CNo$ over $\mathbf{N}^{o}$. In this particular case, one has at least a canonical arrow
\[
C \longrightarrow C_{\infty}
\]
about which one can ask oneself the question whether it is
\msp{page 109}
a Hot-equivalence. If not, this will prove that in the general case, the Hot-type of $C$ (and a fortiori its isomorphism class in Cat\,hot) has no interesting geometric meaning. That is to say that $C$ is interesting only \textit{as a fibered category over $T$}, making it possible in particular to construct the ind-object $H_{T}$ of Cat\,hot or of Hot, but not as a ``bare category''.

But I have just convinced myself that $\CNo \longrightarrow C_{\infty}$ is indeed a Hot-equivalence. So the arrow $\CNo \longrightarrow C$ defines an arrow \textit{in} Hot
\[
C_{\infty} \longrightarrow C
\]
(by composing with the inverse in Hot of the preceding one). Is it an isomorphism? So for this question, one is reduced to the question for $\CNo \longrightarrow C$. But if one takes $C = T$, this already poses the same question for $\mathbf{N}^{o} \longrightarrow T$, hence the question whether $T$ is aspheric, under the sole hypothesis that $T_{\mathrm{o}}$ is left filtering and countable at infinity (on the left). But of course this is false (e.g. if $T$ has only one object, it is not aspheric for all that!).

This therefore convinces me that in the absence
% foot: Version 22 décembre 2001
\printedpage{70}
% head: A. Grothendieck : Dérivateurs, Ch. VII, Édition des Universités de Montpellier II et Paris VII
\msp{page 110}
of an additional draconian hypothesis on $T$, the ``bare'' category $C$ is of no particular interest, apart from its $\pi_{0}$ (which is also given by $C_{\infty}$, and by $H_{C}$). It is not certain, to be sure, that \underline{Ty} does not satisfy these ``draconian conditions''. But since ``bare'' $C$ does not fit into a congenial formalism (e.g. one has no geometric interpretation of $\mathrm{Hom}_{\mathrm{Cat\,hot}}(S,C)$, or of $\mathrm{Hom}_{\mathrm{Hot}}(S,C)$), I think that in what follows I can leave it out of account. Exit a)!

\begin{flushright}
{\footnotesize 5.11.90}
\end{flushright}

All the same, I spent the whole day yesterday trying to understand the relation of $C$ and of its type in Hot, to that of $C_{\infty}$. Indeed, since one finally has an arrow in Hot
\[
C_{\infty} \longrightarrow C,
\]
deduced from the diagram in (Hot)
\[
\begin{tikzcd}[column sep=huge, row sep=huge]
\CNo \arrow[r, "\sim"] \arrow[d] & C_{\infty} \\
C, &
\end{tikzcd}
\]
to prove that $C_{\infty} \longrightarrow C$ is an isomorphism in Hot
\msp{page 111}
seems the simplest, and perhaps the only, way of proving that the image of $C_{\infty}$ in Hot does not depend, up to unique isomorphism, on the choices made to construct $C_{\infty}$. Not to mention that in the motivating case, $C = \Ch_{\Sigma}(X)$ or $\Ch_{\Sigma}(X;a,b)$, $T = \Ty$, it is the category of paths $C$ that gives rise to an impeccable composition law for paths, functorial, associative and all, whereas there is nothing of the sort for $C_{\infty}$. One does not find a composition law
\[
C_{\infty} \times C_{\infty} \longrightarrow C_{\infty}
\]
in Cat, nor even in Cat\,hot or in Hot - or rather, if one does nevertheless manage to construct an arrow $C_{\infty} \times C_{\infty} \longrightarrow C_{\infty}$ in Cat, it is still laboriously by dint of $\pm$ arbitrary choices, notch by notch, and without knowing whether what one obtains has a uniqueness property, even if only in Hot; one finds a tangible geometric interpretation of this composition only in $\Hotz$ (where one has a genuine groupoid law), or in $H_{T}$ as an object of Ind\,Cat\,hot
% foot: Version 22 décembre 2001
\printedpage{71}
% head: A. Grothendieck : Dérivateurs, Ch. VII, Édition des Universités de Montpellier II et Paris VII
\msp{page 112}
or of Ind\,Hot. (So the direct limit in $\Hotz$ is none other than the element of $\Hotz$ defined by $C_{\infty}$.)

My efforts to prove that
\[
C_{\infty} \longrightarrow C
\]
was an isomorphism in Hot, or even just a 1-isomorphism (i.e. an isomorphism not only for the $\pi_{0}$, which is always true, but also for the $\pi_{1}$ of all the components), have failed. I found that it was enough for the functor % [sic] ont échoués -> ont échoué, accord fautif du participe
\[
T \longrightarrow T_{\mathrm{o}}
\]
to be either smooth and coaspheric, or proper and aspheric - so everything is OK if $T = T_{\mathrm{o}}$, i.e. if $T$ is ordered. But in the case $T = \Ty$, the functor $T \longrightarrow T_{\mathrm{o}}$ is neither proper, nor smooth (not even 0-proper or 0-smooth).

Finally I have come to the conviction that in the case we are concerned with ($T = \Ty$, $C = \Ch_{\Sigma}(X;a,b)$), the arrow $C_{\infty} \longrightarrow C$ is in general \textit{not} an isomorphism in Hot. Indeed, in the case $a = b$, if this \texttt{[were true]},
\msp{page 113}
one would have a description of the Hot-group $\Omega(X,a) = \Ch(X,a)$ (\footnote[23]{Here we take $\Sigma = \mathrm{Fl}(X)$.}) in terms of a category $C$ equipped with an \textit{associative} (and unital) composition law \textit{in} Cat, giving back the group law of $\Omega$ on passing from (Cat) to Hot. Taking the nerve of $C$, one would find a semi-simplicial monoid which, in Hot, would give the group law of $\Omega(X,a)$. Now I remember (\footnote[24]{This is an error - it must be for \textit{other} group objects of Hot than loop groups. I saw in Quillen that $\Omega(X,x)$ must be describable by a semi-simplicial group (when $X$ is a semi-simplicial set, $x$ a vertex). So the problem whether $C_{\infty} \simeq C$ in Hot remains completely open.}) that it was explained to me that in general, the group law of $\Omega(X,a)$ could not be obtained by an \textit{associative} law on a semi-simplicial set. Thus, the failure of my efforts seems to be due, not to my clumsiness, but to the nature of things. % [sic] C'est un erreur (note 24) -> C'est une erreur, genre

I am going to summarize the positive facts I have arrived at in the following diagram, which takes place in the general situation $C \longrightarrow T$ considered at the beginning of the present paragraph:

\msp{page 114}
% foot: Version 22 décembre 2001
\printedpage{72}
% head: A. Grothendieck : Dérivateurs, Ch. VII, Édition des Universités de Montpellier II et Paris VII
\[
\begin{tikzcd}[column sep=small, row sep=huge]
 & & C_{t'} \arrow[d, "{u^{*}\quad\lab{(for $u : t \longrightarrow t'$\\ in $T$)}}"] \arrow[ldd, dotted] \arrow[rdd, dotted] & \\
 & & C_{t} \arrow[ld, swap, "{\lab{arrow defined in\\ Ind\,Cat\,hot\\ (\texttt{[or in]}\\ IndHot,\\ Ind\,($\Hotz$))}}"] \arrow[rd, "{\lab[r]{arrow in Cat\\ (hence in\\ Cat\,hot,\\ Hot,\\ $\Hotz$)}}"] & \\
\varinjlim C_{t} = & |[draw, circle]| H_{T} \arrow[rd, swap, "{\lab{arrow\\ defined in\\ Ind\,Cat\,hot\\ or in Ind\,Hot\\ or in Ind\,$\Hotz$}}"] & \lab{commutative\\ diagram\\ in Ind\,Hot\\ or Ind\,($\Hotz$)} & C \\
\lab{in Ind\,Cat\,hot\\ (or in Ind\,Hot\\ or in Ind\,($\Hotz$))} & & |[draw, circle]| C_{\infty} \arrow[ru, swap, "{\lab[r]{arrow\\ defined\\ in Hot,\\ hence in $\Hotz$,\\ in general not an\\ isomorphism, even\\ in $\Hotz$}}"] & \\
 & & = \varinjlim H_{T} \text{ in } \Hotz &
\end{tikzcd}
\]

\texttt{[The annotations in this diagram concern the solid arrows.]}

The important objects are $H_{T}$ and $C_{\infty}$, whereas $C$ (over $T$) and the $C_{t}$ take on an intermediate technical role, to construct $H_{T}$ and $C_{\infty}$. One can consider $C_{\infty}$ as a category (object of (Cat)), but whose construction depends on many arbitrary choices, unlike $H_{T}$ in Ind\,Cat\,hot. It is the object of $\Hotz$ deduced from $C_{\infty}$ that has a canonical meaning, and can be described as the direct limit of $H_{T}$ as an object of Ind\,$\Hotz$, i.e. as representing the same functor on $\mathrm{Hot}_{0}$ as $H_{T}$.

\msp{page 115}
In the geometric case $T = \Ty$, and $C$ of the form $\Ch_{\Sigma}(X;a,b)$, giving rise to $C_{\infty}(a,b)$, one \textit{does not have} a natural definition of a bifunctor of composition of paths
\[
C_{\infty}(a,b) \times C_{\infty}(b,c) \ \longrightarrow\ C_{\infty}(a,c).
\]
One can define, to be sure, pairings
\[
\Ch_{\Sigma,\tau}(a,b) \times \Ch_{\Sigma,\tau'}(b,c) \ \longrightarrow\ \Ch_{\Sigma,\tau\tau'}(a,c),
\]
or, with the notation relating to the datum of
\[
\begin{aligned}
&\underline{t} : \mathbf{N}^{o} \longrightarrow \Ty,\\
&\Ch_{n}(a,b) \times \Ch_{m}(b,c) \longrightarrow \Ch_{n+m}(a,c),
\end{aligned}
\]
% foot: Version 22 décembre 2001
\printedpage{73}
% head: A. Grothendieck : Dérivateurs, Ch. VII, Édition des Universités de Montpellier II et Paris VII
but these pairings \textit{are not} compatible with the transition morphisms of the three direct systems of categories $(\Ch_{n})_{n\in\mathbf{N}}$, they are so \textit{only up to homotopy.} This is enough, however, to have perfect pairings
\[
H_{T}(a,b) \times H_{T}(b,c) \ \longrightarrow\ H_{T}(a,c),
\]
with associativity, and (if $\Sigma = \mathrm{Fl}(X)$ in its entirety) existence of \textit{true} inverses, i.e. of arrows (in Ind\,Cat\,hot) of ``passage to the opposite path''

\msp{page 116}
\[
H_{T}(a,b) \ \longrightarrow\ H_{T}(b,a), \qquad c \longmapsto c^{o}, \tag*{$(*)$}
\]
giving rise to the commutative diagram expressing that this is indeed an inverse for composition. Another way of saying it is that for every object $Y$ of Cat\,hot (or of Ind\,Cat\,hot), the sets
\[
\mathrm{Hom}_{\mathrm{Cat\,hot}}(Y, H_{T}(a,b)) \overset{\text{def}}{=} H_{T}^{Y}(a,b)
\]
form, for variable $a,b \in X$, the arrows of a \textit{groupoid} (\footnote[25]{If $\Sigma = \mathrm{Fl}(X)$. If $\Sigma$ is smaller, instead of a groupoid, one finds a \textit{category}, nothing more.}) $\mathrm{Ch}_{\Sigma}^{Y}(X)$ having as set of vertices $\mathrm{Ob}\,X$, and $H_{T}^{Y}(a,b)$ being made up of the arrows from $a$ to $b$. The map
\[
H_{T}^{Y}(a,b) \longrightarrow H_{T}^{Y}(b,a), \qquad c \longmapsto c^{o},
\]
deduced from the functor $(*)$ above is none other than the map of passage to the inverse. This category is moreover none other than the one deduced from % [sic] celle déduit -> celle déduite, accord
\[
X(Y)\Sigma_{Y}^{-1}, \qquad \text{where } \left\{
\begin{aligned}
&X(Y) = \underline{\mathrm{Hom}}_{\mathrm{Cat}}(Y,X)\\
&\Sigma_{Y} \subseteq \mathrm{Fl}(X(Y)),\\
&\text{defined in terms of } \Sigma \text{ ``argument by argument''},
\end{aligned}
\right.
\]
by means of the functor
\[
X \longrightarrow X(Y), \qquad a \longmapsto a_{Y},
\]
\msp{page 117}
associating to every $a \in X$ the \textit{constant} functor on $Y$ with value $a$, taking
\[
H_{T}(a,b) \simeq \mathrm{Hom}_{X(Y)\Sigma_{Y}^{-1}}(a_{Y}, b_{Y}) \ .
\]

In the same geometric case, I have convinced myself, moreover, that $C_{\infty}$ was, as an object of Hot (not only of $\Hotz$) \textit{unique up to} Hot-\textit{isomorphism,} but (strange thing) perhaps not up to \textit{unique} Hot-isomorphism. In general, returning to the
% foot: Version 22 décembre 2001

\end{document}
